\chapter{Đánh đổi và thảo luận}
\label{ch:discussion}

\section{Những gì kết quả chứng minh}

Các phép đo cho thấy ba điều trong bối cảnh này. Hợp nhất tĩnh gần với dự đoán hằng trên hai tác vụ nhỏ nhất. Các bộ định tuyến huấn luyện trên 256 mẫu mỗi tác vụ nâng điểm trung bình thêm 12.5 điểm và mức tăng so với lớp đa số từ 18.4 lên 59.1. Các hệ số cố định huấn luyện trên cùng dữ liệu chỉ cải thiện điểm trung bình 1 điểm.

Diễn giải của chúng tôi là định tuyến phụ thuộc đầu vào đóng góp phần lớn cải thiện so với Task Arithmetic. Diễn giải này vượt quá phép đo trực tiếp, vì hai biến thể còn khác nhau về sức chứa (7.1M so với 60 tham số) và learning rate, và ablation chỉ được thử với một learning rate, dù điểm dev tốt nhất của nó chỉ thay đổi 0.2.

\section{Đánh đổi giữa độ chính xác, tham số/lưu trữ và độ trễ}

WEMoE full MLP và intermediate-only đạt điểm trung bình tương tự (69.9 và 69.1), nhưng thất bại ở các tác vụ khác nhau. Intermediate-only lưu 47.2\% tham số của năm mô hình, nhỏ hơn biến thể full 35\%, huấn luyện nhanh gấp khoảng hai lần (8.4 so với 16.3 phút mỗi lần chạy) và suy luận nhanh hơn. Full MLP có mức tăng cao hơn (59.1 so với 54.8), dù khoảng cách 0.8 điểm ở điểm trung bình nằm trong mức nhiễu lấy mẫu.

Lưu ít tham số hơn không có nghĩa là suy luận nhanh hơn. Cả hai biến thể WEMoE lưu ít hơn năm chuyên gia nhưng lại chậm hơn trên mỗi mẫu, vì mỗi mẫu có trọng số hợp nhất riêng và không thể dùng chung một phép nhân ma trận trong batch. Task Arithmetic nhỏ nhất (20\%) và có tốc độ như một BERT, nhưng điểm thấp hơn WEMoE full MLP 12.5 điểm.

\section{Ý nghĩa thực tiễn}

Mỗi đầu vào trong thiết lập này đều đi kèm mã tác vụ, nên năm chuyên gia có thể được chọn trực tiếp theo tác vụ. Chúng chính xác hơn WEMoE full MLP trên mọi tác vụ, nhanh hơn trên mỗi mẫu và chỉ lớn hơn khoảng 37\%. Khi đã biết mã tác vụ và dung lượng cho phép, các chuyên gia vẫn là lựa chọn tốt hơn trong thực nghiệm này. WEMoE chủ yếu đáng cân nhắc khi cần giảm lưu trữ, và trường hợp này mạnh nhất với biến thể intermediate-only. Các bối cảnh không có mã tác vụ tường minh có thể có lợi hơn cho nó, nhưng mô hình hợp nhất ở đây vẫn dùng mã tác vụ để chọn head, và các bối cảnh đó chưa được đánh giá. Cũng không nên chỉ đọc điểm trung bình, vì nó che khuất những thất bại ở mức tác vụ mà mốc lớp đa số làm lộ ra.

\section{Hạn chế, nguy cơ đối với tính hợp lệ và câu hỏi còn mở}

\begin{itemize}
  \item \textbf{Tính hợp lệ nội tại.} ``Few-shot'' chưa nghiêm ngặt theo nghĩa của Perez và cộng sự~\cite{perez2021truefewshot}, vì nhãn dùng để lựa chọn cũng phải được tính: ở $k=256$, FSFT huấn luyện trên tổng cộng 1\,280 mẫu có nhãn, trong khi việc lựa chọn dùng 3\,471 mẫu dev. Điểm dev của SST-2 cũng đánh giá cao hơn thực tế chất lượng mô hình hợp nhất (Mục~\ref{sec:merging}).
  \item \textbf{Tính hợp lệ thống kê.} Độ lệch chuẩn bao gồm ba seed FSFT và không tính biến thiên của các chuyên gia và cách chia dữ liệu, vốn chỉ dùng một seed. RTE và MRPC có tập kiểm định nhỏ, với sai số chuẩn khoảng 3 và 2 điểm, lớn hơn một số khác biệt được thảo luận, kể cả 52.0 so với 52.7 trên RTE.
  \item \textbf{Tính hợp lệ bên ngoài.} Kết quả đến từ một lần chạy đầy đủ trên năm tác vụ nhị phân với một backbone, và không nên khái quát ra ngoài thiết lập này. Backbone, tác vụ, hàm mục tiêu và dữ liệu đều khác với Tang và cộng sự~\cite{tang2024wemoe}, nên mọi so sánh với nghiên cứu đó chỉ mang tính định tính.
  \item \textbf{Tính hợp lệ của cấu trúc đo.} Điểm trung bình trộn lẫn accuracy, F1 và MCC. Độ trễ dùng đầu vào ngẫu nhiên trên một GPU, và số liệu batch 1 dao động giữa các lần đo. Các số liệu định tuyến chỉ cho thấy mối liên hệ, và một số trong đó dao động giữa các seed.
\end{itemize}

Các câu hỏi còn mở gồm vì sao RTE thất bại với full MLP, biến thiên định tuyến trong cùng một tác vụ đi theo yếu tố nào, và điểm còn tăng bao xa khi vượt quá $k=256$.

\section{Phát hiện đã chứng minh và giả thuyết còn mở}

Bảng~\ref{tab:findings} tách những gì phép đo xác lập trong bối cảnh này khỏi những gì còn phải kiểm chứng.

\begin{table}[htbp]
  \centering\tablestyle
  \caption{Phát hiện đã chứng minh và giả thuyết còn mở.}
  \label{tab:findings}
  \begin{tabularx}{\linewidth}{*{2}{>{\raggedright\arraybackslash}X}}
    \toprule
    Đã chứng minh trong bối cảnh này & Giả thuyết còn mở\\
    \midrule
    FSFT với bộ định tuyến cải thiện rõ rệt so với Task Arithmetic về trung bình. & Tính phụ thuộc đầu vào, chứ không phải sức chứa, giải thích phần lớn cải thiện này.\\
    Hệ số cố định với các thiết lập đã thử gần như không cải thiện. & Một learning rate khác sẽ không thay đổi kết quả này.\\
    Ngân sách nhỏ không ổn định; điểm trung bình vẫn tăng ở $k=256$. & Ngân sách lớn hơn sẽ tiếp tục nâng điểm.\\
    Hai phạm vi yếu ở các tác vụ khác nhau. & Thất bại ở RTE bắt nguồn từ đánh đổi trong hàm mục tiêu tổng và việc chọn theo điểm trung bình, hoặc từ chuyên gia RTE yếu.\\
    Intermediate-only đạt bước tốt nhất muộn (400--500 trên 500). & Thêm bước sẽ cải thiện mô hình intermediate-only.\\
    \bottomrule
  \end{tabularx}
\end{table}
